\subsection*{The mathematics in one view}
A nonzero lattice point $T=(a,b)$ is \textbf{visible} from $O=(0,0)$ when the open line segment $OT$ contains no other lattice point. A blocker must lie exactly on that segment; O and T are not blockers. The packet uses the first quadrant, including the axes.

\textbf{Exact criterion.} For nonnegative integer coordinates, not both zero, let $d=\gcd(a,b)$. The first lattice point on the ray is $(a/d,b/d)$, and the interior points are exactly $k(a/d,b/d)$, for $k=1,\dots,d-1$. Thus T is visible if and only if $d=1$, and there are $d-1$ blockers. On an axis use $\gcd(a,0)=a$; only the first unit point is visible.

\textbf{What the geometry encodes.} Every lattice ray has one primitive first point, with coprime coordinates. Its other lattice points are positive integer multiples of that point. Scaling a direction explains hidden dots; a common divisor measures how many equal primitive steps reach the target.

\textbf{Limits.} These are ideal points, not thick trees: an almost-collinear dot does not block. On these rectangular grids every intermediate lattice point is present. Visibility could differ if dots were removed. The source's asymptotic density discussion is not part of this lesson.

\subsection*{What each group is doing}
\textbf{K--1:} Test exact sight lines, find first blockers, and compare rows and rays. Counting to 6 and holding a thread straight suffice; coordinates may be spoken or pointed to. A demonstration of a blocker is a valid explanation of a particular hidden target.

\textbf{Grades 2--3:} Group dots sharing a direction and conjecture a divisibility rule from examples, then explain a fully visible row. Coordinates mean ``across, up''; small equal groups and multiplication are helpful.

\textbf{Grades 4--5:} Explain both directions of the visibility criterion and why the blocker list is complete. Division, common factors and proportional steps are the prerequisites. The reduced-fraction argument can stay with the facilitator if children can only justify particular rays.
